\documentclass[12pt]{article}
\usepackage[utf8]{inputenc}
\usepackage[T1]{fontenc}
\usepackage{hyperref}
\title{Yardliyawarra Language Resources for K-12 Curriculum Development}
\author{Compiled by Claude Sonnet 5.5 (claude-sonnet-5-5) for Andr\'as Kornai}
\date{}
\begin{document}
\maketitle

\section{Introduction}
Yardliyawarra (ISO 639-3: yxl; AIATSIS spelling Yardliyawara; also Jadliaura, Yadliyawarra) is the language of the Yardliyawara people of the northern Flinders Ranges, South Australia. It was spoken on the eastern side of the northern Flinders Ranges. Its classification is contested: it has been placed with the Karnic languages, and Hercus and Austin group it in the Yarli subgroup. The Yardliyawara are now counted within the wider Adnyamathanha people. I found no information on current speakers and no dedicated teaching materials.

\section{Primary Educational Resources}
I found no primers or textbooks specific to Yardliyawarra. Adnyamathanha, which is closely linked and still has speakers, has more resources; its first bilingual dictionary appeared in November 2020.

\section{Linguistic and Reference Materials}
\begin{itemize}
\item Wikipedia, Yardliyawara: \url{https://en.wikipedia.org/wiki/Yardliyawara}
\item AIATSIS Austlang record for the language: \url{https://collection.aiatsis.gov.au/austlang/language/L7}
\item Wikipedia, Adnyamathanha language: \url{https://en.wikipedia.org/wiki/Adnyamathanha_language}
\end{itemize}

\section{Religious and Community Resources}
I found no scripture or community organisation specific to Yardliyawarra.

\section{Digital and Audio Resources}
\begin{itemize}
\item Mobile Language Team (South Australia), a language programme that supports Aboriginal languages in the region; I did not verify its coverage of Yardliyawarra: \url{https://mobilelanguageteam.com.au/?p=229}
\item The Yarli Project, a Corner Country document that appeared in search results: \url{https://visitcornercountry.com.au/wp-content/uploads/2024/03/The-Yarli-Project-4.pdf}
\end{itemize}

\section{Curriculum Development Context}
Material specific to Yardliyawarra is sparse. A programme would probably rely on Adnyamathanha resources and Yarli-region documentation, developed with Adnyamathanha community members.
\end{document}